\documentclass[11pt,a4paper]{article}
\PassOptionsToPackage{pdfusetitle=false}{hyperref}
\usepackage[utf8]{inputenc}
\usepackage[T1]{fontenc}
\usepackage[margin=2.5cm]{geometry}
\usepackage{amsmath,amssymb}
\usepackage{booktabs,array,longtable}
\usepackage{xcolor}
\usepackage{tcolorbox}
\tcbuselibrary{skins,breakable}
\usepackage{enumitem}
\usepackage[pdfusetitle=false]{hyperref}
\usepackage[numbers]{natbib}
\bibliographystyle{unsrtnat}
\hypersetup{colorlinks=true,linkcolor=blue!60!black,urlcolor=blue!60!black,
            citecolor=blue!60!black,
            pdftitle={Falsifiability: How to Test the GTE Framework},
            pdfauthor={Nova Spivack}}
\setlength{\emergencystretch}{2em}
\usepackage{tikz}
\usetikzlibrary{arrows.meta,positioning,shapes.geometric,calc,fit,backgrounds}

%─── Color palette ───────────────────────────────────────────────────────────
\definecolor{boxblue}{RGB}{220,235,255}
\definecolor{boxorange}{RGB}{255,240,210}
\definecolor{boxgreen}{RGB}{220,245,220}
\definecolor{boxred}{RGB}{255,230,225}
\definecolor{boxpurple}{RGB}{240,225,255}
\definecolor{boxyellow}{RGB}{255,252,210}
\definecolor{boxteal}{RGB}{210,240,238}

%─── Macros ──────────────────────────────────────────────────────────────────
\newcommand{\term}[1]{\textbf{#1}}
\newcommand{\lean}[1]{%
  \penalty-100
  \begingroup\def\_{\textunderscore\allowbreak}\texttt{#1}\endgroup}
\newcommand{\CatAL}{\textbf{CatAL}}
\newcommand{\CatAD}{\textbf{CatAD}}
\newcommand{\CatA}{\textbf{CatA}}
\newcommand{\PHIMDL}{$\Phi_{\rm MDL}$}
\newcommand{\sigW}{\sin^2\!\theta_W}

\newcommand{\TutorialDOI}{10.5281/zenodo.20669142}
\date{2026}

\title{\textbf{Falsifiability: How to Test the GTE Framework}\\[6pt]
\large The Complete Prediction and Falsification Register\\[8pt]
\normalsize
\href{https://novaspivack.com/research}{novaspivack.com/research}\texorpdfstring{\quad}{ }
\href{https://doi.org/10.5281/zenodo.20560550}{P48 complete theory monograph}\\[4pt]
\small\href{https://doi.org/\TutorialDOI}{doi.org/\TutorialDOI}
  \textit{(registration pending)}\\[4pt]
\normalsize\textit{UGP Physics Tutorial Series}}
\author{Nova Spivack}

\begin{document}
\setcounter{tocdepth}{2}
\maketitle
\tableofcontents
\newpage

\begin{tcolorbox}[colback=blue!6,colframe=blue!40!black,
  title=\textbf{UGP Physics Tutorial Series}]
\textbf{Prerequisites (read first):}
  \href{https://doi.org/10.5281/zenodo.20657889}%
    {\emph{The GTE Polynomial: A Step-by-Step Cheat Sheet}}
  $\cdot$
  \href{https://doi.org/10.5281/zenodo.20657895}%
    {\emph{Perfect Self-Containment and MDL}}\\[4pt]
\textbf{Next tutorials:}
  None --- this is a terminal reference tutorial.\\[4pt]
\textbf{See also:}
  \href{https://doi.org/10.5281/zenodo.20657913}%
    {\emph{From the Arithmetic Substrate to Particle Masses}}
  $\cdot$
  \href{https://doi.org/10.5281/zenodo.20657923}%
    {\emph{The Weinberg Angle from $\mathbb{Z}_7$ Arithmetic}}
  $\cdot$
  \href{https://doi.org/10.5281/zenodo.20657861}%
    {\emph{The Born Rule as a Theorem}}
  $\cdot$
  \href{https://doi.org/10.5281/zenodo.20669152}%
    {\emph{New Physics Predictions: What GTE Predicts Will and Won't Be Found}}
\end{tcolorbox}


%══════════════════════════════════════════════════════════════════════════════
\section{Why Falsifiability Matters for GTE}
\label{sec:intro}
%══════════════════════════════════════════════════════════════════════════════

\begin{tcolorbox}[colback=boxblue,colframe=blue!40!black,
  title=The core argument]
A theory without free parameters is maximally falsifiable.  GTE has zero
continuously adjustable parameters: every numerical output is a derivation, not
a fit.  This makes GTE \emph{more} falsifiable than the Standard Model, which
has 19 free parameters (plus at least 3 for neutrinos) that can absorb
discrepancies.  A single confirmed measurement outside a GTE prediction band,
at the stated precision, would refute the framework with no available escape
route.
\end{tcolorbox}

\term{GTE} (\term{Generative Triple Evolution}) derives all Standard Model
parameters from a single 19-bit algebraic certificate: the polynomial
$p(L,C,R) = C + R - CR - LCR$ over $\mathrm{GF}(7)$, selected by the
\term{MDL} (\term{Minimum Description Length}) principle.  With zero free
parameters, every output of the framework is a prediction that could in principle
be falsified.

This tutorial is the experimentalist's guide to GTE.  It presents the complete
prediction register organized by sector, identifies the strongest confirmations,
explains the current tensions honestly, provides clean falsification criteria for
each sector, and shows how each experimental community can contribute to testing
the framework.

All comparisons use PDG~2024 \cite{SpivackSM} and NuFIT~6.0 IC24 NH
\cite{SpivackGTECompleteFramework} for neutrino parameters.

%══════════════════════════════════════════════════════════════════════════════
\section{The Complete Prediction Register}
\label{sec:register}
%══════════════════════════════════════════════════════════════════════════════

\subsection{Sector 1: Particle Masses}

Nine charged-fermion masses spanning six orders of magnitude — three leptons and
six quarks — are derived from the GTE arithmetic cascade with zero free
parameters \cite{SpivackSM, SpivackGTECompleteFramework}.

\bigskip\medskip
\begin{table}[htbp]
\centering
\caption{Charged-fermion mass predictions vs.\ PDG~2024.  All GTE predictions
are derived from the InformationMassTransformer pipeline with the Elegant Kernel.
The tau mass is the single anchor fixed by the Self-Consistency Condition; all
other masses are outputs.  RMS = $0.295\%$ across all nine.}
\label{tab:masses}
\setlength{\tabcolsep}{4pt}
\begin{tabular}{lrrccc}
\toprule
Fermion & PDG 2024 & GTE prediction & $\sigma$ & Cert. & Falsification threshold \\
\midrule
$e$ & $0.5110$~MeV & $0.511$~MeV & $0.00$ & \CatAL & --- (anchor) \\
$\mu$ & $105.66$~MeV & $105.66$~MeV & $0.00$ & \CatAL & --- (anchor) \\
$\tau$ & $1776.86$~MeV & $1776.86$~MeV & $0.00$ & anchor & --- \\
$u$ & $2.16$~MeV & $2.23$~MeV & $<0.1$ & \CatA & outside $[2.0,2.5]$~MeV \\
$d$ & $4.67$~MeV & $4.55$~MeV & $<0.2$ & \CatA & outside $[4.0,5.5]$~MeV \\
$s$ & $93.4$~MeV & $93.0$~MeV & $<0.1$ & \CatA & outside $[85,100]$~MeV \\
$c$ & $1275$~MeV & $1277$~MeV & $<0.1$ & \CatA & outside $[1200,1350]$~MeV \\
$b$ & $4183$~MeV & $4184$~MeV & $<0.04$ & \CatAD & outside $[4100,4300]$~MeV \\
$t$ & $172570$~MeV & $172610$~MeV & $<0.03$ & \CatA & outside $[170,175]$~GeV \\
\midrule
Pion $\pi^\pm$ & $139.5703$~MeV & $139.57$~MeV & $0.00$ & \CatAL & outside $[139.0,140.1]$~MeV \\
\bottomrule
\end{tabular}
\end{table}

The \textbf{Koide relation} $Q = (m_e + m_\mu + m_\tau)/(\sqrt{m_e}+\sqrt{m_\mu}+\sqrt{m_\tau})^2 = 2/3$
is a theorem (\lean{ucl\_lepton\_sector\_koide\_identity}, \CatAL), not a fit.

\subsection{Sector 2: Mixing Angles and CP Violation}

\begin{table}[htbp]
\centering
\caption{Mixing-sector predictions vs.\ PDG~2024 and NuFIT~6.0 IC24 NH.}
\label{tab:mixing}
\setlength{\tabcolsep}{4pt}
\begin{tabular}{lllcc}
\toprule
Observable & GTE prediction & PDG / NuFIT & $\sigma$ & Cert. \\
\midrule
Weinberg angle $\sigW$ & $0.23129$ (two-loop) & $0.23121 \pm 0.00004$ & $+0.038$ & \CatAD \\
CKM CP phase $\delta_{\rm CKM}$ & $68.51^\circ$ & $68.5^\circ \pm 2.5^\circ$ & $0.00$ & \CatAL \\
PMNS CP phase $\delta_{\rm CP}$ & $4/7 \times 360^\circ = 205.71^\circ$ & $212^\circ \pm 26^\circ$ & $-0.15$ & \CatA \\
PMNS solar $\sin^2\!\theta_{12}$ & $4/13 = 0.3077$ & $0.308 \pm 0.012$ & $-0.02$ & \CatAL \\
PMNS atm $\sin^2\!\theta_{23}$ & $19/42 = 0.4524$ & $0.470 \pm 0.015$ & $-1.0$ & \CatAL \\
PMNS reactor $\sin\theta_{13}$ & $11/73 = 0.1507$ & $0.1489 \pm 0.004$ & $+0.23$ & \CatAL \\
\bottomrule
\end{tabular}
\end{table}

All three PMNS angles are derived from orbit-ratio formulas with zero free
parameters (\lean{pmns\_solar\_sin\_sq}, \lean{pmns\_atm\_sin\_sq},
\lean{pmns\_reactor\_sin\_val}, all \CatAL) \cite{SpivackGTECompleteFramework}.

\subsection{Sector 3: Cosmology}

\begin{table}[htbp]
\centering
\caption{Cosmological predictions vs.\ Planck~2018 and other measurements.}
\label{tab:cosmo}
\setlength{\tabcolsep}{4pt}
\begin{tabular}{llllc}
\toprule
Observable & GTE prediction & Measurement & $\sigma$ & Cert. \\
\midrule
CMB spectral tilt $n_s$ & $1 - \frac{\ln 2}{2\pi^2} = 0.96488$ & $0.9649 \pm 0.0044$ & $-0.005$ & \CatAL \\
Baryon asymmetry $\eta_B$ & $6.109 \times 10^{-10}$ & $6.10 \times 10^{-10}$ & $+0.15$ & \CatAL \\
Dark matter relic $\Omega_{\rm DM} h^2$ & $0.11994$ & $0.12011 \pm 0.0012$ & $-0.15\%$ & \CatAD \\
$\Omega_\Lambda$ (bracket) & $[3\pi/14,\; 0.6899]$ & $0.6889 \pm 0.0056$ & in bracket & \CatAD \\
Primordial tensor ratio $r$ & $= 0$ exactly & $< 0.036$ & consistent & \CatA \\
Dark energy $w$ & $= -1$ exactly & $-1.006 \pm 0.045$ & consistent & \CatAD \\
Neutrino mass sum $\Sigma m_\nu$ & $59.4$~meV & $< 120$~meV & consistent & \CatA \\
Normal mass ordering & NH only & NH preferred & consistent & \CatAD \\
\bottomrule
\end{tabular}
\end{table}

The spectral tilt $n_s = 0.96488$ is derived from GTE's $\mathbb{Z}_7$ defect
cosmology with zero parameters, machine-certified
(\lean{cmb\_spectral\_index\_equals\_gte\_prediction}, \CatAL)
\cite{SpivackGTECosmology}.  The baryon asymmetry $\eta_B = 6.109 \times 10^{-10}$
is derived via the FKTT (Fermi-Kink Top-Top) baryogenesis mechanism, machine-certified
(\lean{kink\_top\_coupling\_eq\_eps\_FN}, \lean{phi\_mdl\_kink\_bps\_saturation},
\CatAL) \cite{SpivackGTECosmology}.

\subsection{Sector 4: Constants and Structure}

\begin{table}[htbp]
\centering
\caption{Structural and constant predictions.}
\label{tab:constants}
\setlength{\tabcolsep}{4pt}
\begin{tabular}{lllcc}
\toprule
Observable & GTE prediction & Value / status & $\sigma$ & Cert. \\
\midrule
Fine-structure constant $\alpha^{-1}$ & $2^7 + 3^2 = 137$ & $137.036$ & $\sim 0.026\%$ & \CatA \\
Newton's constant $G_N$ & $m_\tau^2/M_{\rm Pl}^2$ & $0.040\%$ from PDG & $\sim 0.04$ & \CatAD \\
Strong CP angle $\theta_{\rm QCD}$ & $= 0$ exactly ($F_{21}$) & $< 10^{-10}$ & consistent & \CatAL \\
Spacetime dimensions & $3+1$ (forced) & $3+1$ observed & exact & \CatAD \\
EFT threshold $\Lambda_{\rm GTE}$ & $\approx 2.03$~GeV & near hadronic scale & target & \CatAD \\
\bottomrule
\end{tabular}
\end{table}

\subsection{Sector 5: Dark Sector and New Physics}

\begin{table}[htbp]
\centering
\caption{Dark sector and new physics predictions.}
\label{tab:dark}
\setlength{\tabcolsep}{4pt}
\begin{tabular}{lllcc}
\toprule
Observable & GTE prediction & Status & Cert. & Experiment \\
\midrule
Dark lepton GTE-P7 & $211.9$~MeV & unobserved (predicted) & \CatAD & Belle~II \\
Dark matter $\chi_1$ & $0.54$~MeV UFIDM & $\sigma_{\rm SI} \sim 10^{-61}$~cm$^2$ & \CatAD & --- \\
Dark photon $A'$ & \textbf{does not exist} & null searches confirmed & Lean & NA64, Belle~II \\
Dim-4 proton decay & \textbf{forbidden} & $\tau_p > 1.6 \times 10^{34}$~yr & \CatAL & Hyper-K, DUNE \\
Axion / ALP & \textbf{does not exist} & null searches confirmed & \CatAL & ADMX, CASPEr \\
\bottomrule
\end{tabular}
\end{table}

%══════════════════════════════════════════════════════════════════════════════
\section{The Strongest Confirmations}
\label{sec:confirmations}
%══════════════════════════════════════════════════════════════════════════════

Among the 40+ quantitative GTE predictions, the following are the most
precisely confirmed — predictions the framework was required to satisfy and did.


%══════════════════════════════════════════════════════════════════════════════
\section{Current Tensions: An Honest Assessment}
\label{sec:tensions}
%══════════════════════════════════════════════════════════════════════════════

Scientific honesty requires presenting tensions clearly, not minimizing them.

\begin{tcolorbox}[colback=boxyellow,colframe=orange!60!black,
  title=Current tensions in the GTE prediction register, breakable]
\begin{itemize}[itemsep=4pt]
  \item \textbf{PMNS atmospheric mixing angle $\theta_{23}$:}
    GTE predicts $\sin^2\!\theta_{23} = 19/42 = 0.4524$ (\CatAL).
    NuFIT~6.0 IC24 NH gives $0.470 \pm 0.015$, which is $-1.0\sigma$ from
    the GTE value.  This is within $1\sigma$ and therefore not a confirmed
    discrepancy, but it is the most notable tension in the mixing sector.
    DUNE and Hyper-Kamiokande will resolve this to $<0.5\sigma$ precision.

  \item \textbf{Cosmological constant (Route 2):}
    The GTE $\Omega_\Lambda$ upper bracket bound $0.6899$ is a bound, not a
    standalone prediction.  The Planck~2018 value $0.6889 \pm 0.0056$ lies
    $0.18\sigma$ below this bound.  The lower bound $3\pi/14 \approx 0.6731$
    is satisfied with $\sim 2.8\sigma$ margin.  The bracket is \CatAD but
    is presented as a bracket, not a point prediction; the theory does not yet
    predict the exact value of $\Omega_\Lambda$ within the bracket.

  \item \textbf{Hubble tension:}
    GTE predicts $H_0 = 67.95$~km/s/Mpc (\CatA), consistent with Planck~2018
    ($67.4 \pm 0.5$~km/s/Mpc, $+1.1\sigma$) but in tension with local distance
    ladder measurements ($73.0 \pm 1.0$~km/s/Mpc, $\sim 5\sigma$).  GTE does
    not resolve the Hubble tension — it predicts a Planck-consistent value and
    does not have a mechanism for a higher $H_0$.

  \item \textbf{Strong coupling $\alpha_s(\Lambda_{\rm GTE})$:}
    GTE predicts $e^2(\Lambda_{\rm GTE}) = 7/2$ (physically normalized),
    which matches PDG $\alpha_s$ to within $+7.4\%$ at one loop (CatAD
    consistency).  The residual offset after the precision kink form-factor
    measurement ($+1.3$--$1.5\sigma$) is assigned to two-loop matching effects.
    An exact-grade upgrade requires the two-loop matching calculation.
\end{itemize}
\end{tcolorbox}

None of these tensions constitutes a confirmed falsification.  Each has a named
pathway to resolution (a specific calculation or measurement).

%══════════════════════════════════════════════════════════════════════════════
\section{Clean Falsification Criteria}
\label{sec:falsification}
%══════════════════════════════════════════════════════════════════════════════

The following is the definitive list of measurements that would falsify GTE at
$>3\sigma$ confidence, with no escape route, assuming the experiment is
correctly calibrated.

\begin{longtable}[l]{@{}p{3.5cm}p{3.2cm}p{2.4cm}p{2.4cm}p{2.5cm}@{}}
\caption{Clean falsification criteria: measurements that would definitively
refute GTE as currently formulated.  Each row gives the GTE prediction, the
measurement precision threshold required for a $>3\sigma$ refutation, the primary
experiment, and the expected timeline.}
\label{tab:falsify}\\
\hline
\textbf{Observable} & \textbf{GTE prediction} &
\textbf{Falsifying measurement} & \textbf{Experiment} & \textbf{Timeline} \\
\hline
\endfirsthead
\hline
\textbf{Observable} & \textbf{GTE prediction} &
\textbf{Falsifying measurement} & \textbf{Experiment} & \textbf{Timeline} \\
\hline
\endhead
\hline
\endfoot
CMB spectral tilt $n_s$ &
  $0.96488$ &
  $n_s > 0.98$ or $<0.95$ at $>3\sigma$ &
  CMB-S4, Simons Obs. &
  2028+ \\[4pt]
Baryon asymmetry $\eta_B$ &
  $6.109 \times 10^{-10}$ &
  outside $[5.9,\,6.4]\times10^{-10}$ at $>3\sigma$ &
  CMB-S4, BBN &
  2028+ \\[4pt]
CKM CP phase $\delta_{\rm CKM}$ &
  $68.51^\circ$ &
  outside $[60^\circ,\,77^\circ]$ at $>3\sigma$ &
  Belle~II, LHCb &
  2027+ \\[4pt]
Weinberg angle $\sigW$ &
  $0.23129$ &
  outside $[0.228,\,0.234]$ at $>3\sigma$ &
  HL-LHC, FCC-ee &
  2035+ \\[4pt]
Strong CP angle $\theta_{\rm QCD}$ &
  $= 0$ exactly &
  any nEDM establishing $\theta \neq 0$ at $>3\sigma$ &
  nEDM, ILL &
  ongoing \\[4pt]
PMNS $\delta_{\rm CP}^{\rm PMNS}$ &
  $205.71^\circ$ &
  outside $[175^\circ,\,240^\circ]$ at $>3\sigma$ &
  DUNE, Hyper-K &
  2030+ \\[4pt]
$\Delta m^2_{21}$ (solar) &
  $7.37\times10^{-5}$~eV$^2$ &
  outside $[7.30,\,7.44]\times10^{-5}$ at $>3\sigma$ &
  JUNO &
  2025+ \\[4pt]
Neutrino mass sum &
  $59.4$~meV &
  $<40$~meV or $>80$~meV at $>3\sigma$ &
  CMB-S4, Euclid &
  2030+ \\[4pt]
GTE-P7 at 211.9~MeV &
  resonance exists &
  no neutral state in $[182,\,243]$~MeV (the pre-registered $\pm14\%$ window) in monophoton &
  Belle~II &
  2031 \\[4pt]
Proton decay (dim-4) &
  forbidden &
  any verified dim-4 decay event &
  Hyper-K, DUNE &
  2035+ \\[4pt]
Dark photon $A'$ &
  does not exist &
  any dark photon discovery &
  NA64, Belle~II &
  ongoing \\[4pt]
Axion / ALP &
  does not exist &
  any confirmed axion signal &
  ADMX, CASPEr &
  ongoing \\[4pt]
Primordial tensors $r$ &
  $r = 0$ exactly &
  $r > 0.01$ at $>3\sigma$ &
  CMB-S4, LiteBIRD &
  2028+ \\[4pt]
Inverted neutrino hierarchy &
  NH only &
  definitive IH determination &
  DUNE, Hyper-K &
  2030+ \\
\end{longtable}

%══════════════════════════════════════════════════════════════════════════════
\section{How to Use This Tutorial: A Guide for Each Community}
\label{sec:communities}
%══════════════════════════════════════════════════════════════════════════════

\subsection{For CMB Observers}

The GTE predictions with the most immediate relevance to CMB science are:
$n_s = 0.96488$ (already $-0.005\sigma$ from Planck; CMB-S4 will test to
$\sigma(n_s) \approx 0.003$), $r = 0$ (LiteBIRD reaches $\sigma(r) \approx 0.001$),
and $\eta_B = 6.109 \times 10^{-10}$ (CMB-S4 will reduce the $\eta_B$ uncertainty
by a factor of $\sim 3$, making the $+0.15\sigma$ a direct comparison).

The most decisive near-term CMB test is $n_s$: a confirmed $n_s > 0.975$ or
$n_s < 0.955$ at $>3\sigma$ from CMB-S4 would refute GTE's defect-cosmology
mechanism.

\subsection{For Collider Physicists}

The primary collider targets are:
\begin{itemize}[itemsep=3pt]
  \item \textbf{Belle~II:} GTE-P7 at 211.9~MeV (highest-priority new-physics
    signal), the dark-photon null (any discovery falsifies), and the dark lepton
    $\chi_2$ at 24.5~MeV.
  \item \textbf{Hyper-Kamiokande / DUNE:} Proton stability (dimension-4 decay
    forbidden; dimension-6 decay via standard GUT mechanism not forbidden).
  \item \textbf{LHC / HL-LHC:} Dark lepton $\chi_3$ at 3.60~GeV (mono-jet or
    Higgs invisible); the EFT threshold near $\Lambda_{\rm GTE} \approx 2$~GeV.
\end{itemize}

\subsection{For Dark Matter Hunters}

GTE predicts UFIDM (ultra-faint interacting dark matter) at $m_{\chi_1} = 0.54$~MeV
with $\sigma_{\rm SI} \sim 10^{-61}$~cm$^2$ — invisible to WIMP-style direct
detection.  Any positive signal from WIMP searches (XENONnT, LZ, PandaX) at
any mass above the neutrino floor would falsify the GTE dark-sector structure
if the signal has the WIMP cross-section profile.

The GTE dark matter candidate can in principle be detected through its
SU(3)$_{\rm dark}$ interactions via stellar cooling constraints and MeV-scale
astrophysical searches.

\subsection{For Lattice QCD}

GTE's pion mass prediction $m_\pi = 139.57$~MeV and the GOR formula as a
theorem are directly testable by lattice QCD.  The prediction uses $m_u = 2.23$~MeV
and $m_d = 4.55$~MeV as inputs; improved lattice determinations of the light quark
masses will sharpen the comparison.

The kink cell count target $\xi_{\rm kink} = 7$ cells (from the physical-point
selection $a = \hbar c/\Lambda_{\rm GTE}$) is a substrate Monte Carlo target:
a lattice calculation of the kink correlation length at the GTE-derived lattice
spacing should find $\xi_{\rm kink} = 7$ (band $6.78$--$7.23 \pm 0.54$).

\subsection{For Lean Formalization}

The GTE Lean~4 library (\href{https://github.com/novaspivack/ugp-lean}{ugp-lean})
currently has 308 canonical modules with zero sorry, covering the key derivations.
Open Lean targets that would upgrade predictions from \CatA\ to \CatAL:

\begin{itemize}[itemsep=3pt]
  \item The sech-overlap integral bound: upgrade $\eta_B$ to full \CatAL
    (currently conditional on two \CatA\ axioms for $I(5)$ and $I(11)$).
  \item The $\lambda_1$-simplicity part of the Perron-Frobenius theorem for the
    spin-7 transfer matrix (awaiting upstream Mathlib support).
  \item The PMNS CP phase $\delta_{\rm CP}^{\rm PMNS} = 4/7 \times 360^\circ$
    (currently \CatA; orbital derivation exists, Lean formalization pending).
\end{itemize}

%══════════════════════════════════════════════════════════════════════════════
\section{The Physical Incompleteness Theorem}
\label{sec:pit}
%══════════════════════════════════════════════════════════════════════════════

\begin{tcolorbox}[colback=boxred,colframe=red!60!black,
  title={Physical Incompleteness Theorem (\CatAD)}]
GTE proves that the exact value of the cosmological constant $\Omega_\Lambda$
is a non-computable real number of Turing degree $0'$ (a member of
$\Delta_2^0 \setminus \Delta_1^0$).  No computable algorithm can determine the
exact realized value; only computable bounds are accessible.\\[4pt]
This means no experiment, however precise, can fully determine $\Omega_\Lambda$
within the GTE framework.  The two bracket endpoints ($3\pi/14$ and $0.6899$)
are computable bounds on the non-computable realized value.  This is analogous
to G\"odel's incompleteness theorem applied to cosmological measurement.
\end{tcolorbox}

This is arguably the single most transformative experimental probe of the GTE
framework: if the Physical Incompleteness Theorem is correct, precision
cosmological measurements will converge on $\Omega_\Lambda$ within the bracket
without ever closing it to a point.  Observations that consistently find
$\Omega_\Lambda \in [0.673, 0.690]$ across independent methods would support this.

%══════════════════════════════════════════════════════════════════════════════
\section{Summary: The GTE Falsifiability Profile}
\label{sec:summary}
%══════════════════════════════════════════════════════════════════════════════

\begin{tcolorbox}[colback=boxblue,colframe=blue!40!black,
  title=Key facts about GTE falsifiability, breakable]
\begin{itemize}[itemsep=4pt]
  \item GTE has \textbf{zero free parameters}: every numerical prediction is
    derived, not fitted.

  \item \textbf{40+ quantitative predictions} have been compared to PDG~2024,
    NuFIT~6.0, and Planck~2018.  Approximately 37 agree within $1\sigma$ or
    match exactly; the remainder have named resolution pathways.

  \item The \textbf{strongest near-term tests} are:
    Belle~II (GTE-P7 at 211.9~MeV), JUNO ($\Delta m^2_{21}$), CMB-S4 ($n_s$
    and $r$), and Hyper-Kamiokande (proton stability).

  \item The \textbf{null predictions} (no axion, no dark photon, no dim-4
    proton decay, no SUSY particles) are confirmed by current searches and
    provide ongoing positive confirmation of the framework.

  \item Every new precision measurement from any frontier experiment is
    simultaneously a test of GTE, because the framework has no parameter room to
    absorb discrepancies.
\end{itemize}
\end{tcolorbox}

\section*{Further Reading}
\begin{tcolorbox}[colback=orange!8,colframe=orange!50!black,
  title=\textbf{Primary Sources}]
The results in this tutorial are derived in the following papers,
all available open-access on Zenodo and at
\href{https://novaspivack.com/research}{novaspivack.com/research}:
\begin{itemize}[itemsep=2pt]
  \item \textbf{P48 --- The Complete GTE Framework} (main synthesis monograph;
    Appendix contains the complete predictions and falsification table):\\
        \href{https://doi.org/10.5281/zenodo.20560550}{doi.org/10.5281/zenodo.20560550}
  \item \textbf{P47 --- Cosmological Predictions of the GTE/$\Phi_{\rm MDL}$ Framework}
        (CMB, baryon asymmetry, cosmological constant):\\
        \href{https://doi.org/10.5281/zenodo.20465803}{doi.org/10.5281/zenodo.20465803}
  \item \textbf{P01 --- Standard Model Parameters from the UGP}
        (all nine charged-fermion masses and the mixing sector):\\
        \href{https://doi.org/10.5281/zenodo.20168787}{doi.org/10.5281/zenodo.20168787}
  \item \textbf{P50 --- The Spin-7 Lattice Model}
        (the kink cell count substrate target and spectral analysis):\\
        \href{https://doi.org/10.5281/zenodo.20603066}{doi.org/10.5281/zenodo.20603066}
\end{itemize}
Lean~4 machine proofs:
\href{https://github.com/novaspivack/ugp-lean}{github.com/novaspivack/ugp-lean}
\end{tcolorbox}

\bibliography{../../papers/bib/Spivack_Papers_Bibliography}
\end{document}
